

\newpage

\subsection*{4. Cálculo del Volumen}

Planteamos la integral triple. Debido a que los límites son constantes, podemos separar la expresión en el producto de tres integrales simples: \textcolor{red}{Primero $W_s$. Habria que sacar el $2$, no?} \textcolor{red}{Escribirlo exactamente igual que ejercicio 25 y luego si separar las integrales. }

\begin{equation*} 
    \text{Vol}(W) =  \iiint =2 \int_{0}^{2\pi} d\theta \int_{0}^{\pi/4} \sin \phi \, d\phi \int_{0}^{1} r^2 \, dr
\end{equation*}

Calculamos cada parte de forma independiente:

\begin{itemize}
    \item \textbf{Integral respecto a $r$:} 
    \[ \int_{0}^{1} r^2 \, dr = \left[ \frac{r^3}{3} \right]_{0}^{1} = \frac{1}{3} \]

    \item \textbf{Integral respecto a $\phi$:} 
    \[ \int_{0}^{\pi/4} \sin \phi \, d\phi = \left[ -\cos \phi \right]_{0}^{\pi/4} = -\frac{\sqrt{2}}{2} - (-1) = 1 - \frac{\sqrt{2}}{2} \]

    \item \textbf{Integral respecto a $\theta$:} 
    \[ \int_{0}^{2\pi} d\theta = \left[ \theta \right]_{0}^{2\pi} = 2\pi \]
\end{itemize}

\vspace{15pt}
\textbf{Resultado final:} \textcolor{red}{Aca $Vol(W) = 2Vol(W_s)$. }
Combinando los resultados parciales y aplicando la simetría multiplicando por 2:

\begin{align*}
    \text{Vol}(W) &= 2 \cdot (2\pi) \cdot \left( 1 - \frac{\sqrt{2}}{2} \right) \cdot \left( \frac{1}{3} \right) \\[10pt]
    \text{Vol}(W) &= \frac{4\pi}{3} \left( 1 - \frac{\sqrt{2}}{2} \right)
\end{align*}
